\documentclass[11pt,a4paper]{article}

% ---------- Encoding & Sprache ----------
\usepackage[utf8]{inputenc}
\usepackage[T1]{fontenc}
\usepackage[ngerman]{babel}
\usepackage{lmodern}

% ---------- Layout ----------
\usepackage[margin=2cm]{geometry}
\usepackage{enumitem}
\usepackage{booktabs}
\usepackage{graphicx}
\usepackage{xcolor}
\usepackage{listings}
\usepackage{microtype}
\usepackage[hidelinks]{hyperref}
\setlist{nosep, itemsep=2pt, topsep=3pt}
\setlength{\parindent}{0pt}
\setlength{\parskip}{5pt}
\setlength{\emergencystretch}{2.5em}

% automatisch erzeugt von install_db.py -- nicht von Hand bearbeiten
\newcommand{\mdTabellen}{17}
\newcommand{\mdDatensaetze}{207}
\newcommand{\mdMinZeilen}{10}
\newcommand{\mdIndizes}{11}
\newcommand{\mdAutoIndizes}{11}
\newcommand{\mdTrigger}{2}
\newcommand{\mdViews}{1}
\newcommand{\mdSeiten}{45}
\newcommand{\mdSeitengroesse}{4096}
\newcommand{\mdBytes}{184.320}
\newcommand{\mdKiB}{180,0}
\newcommand{\mdSQLite}{3.50.4}
\newcommand{\mdZeilenTabelle}{%
\texttt{nutzer} & 12 \\
\texttt{adresse} & 12 \\
\texttt{verlag} & 10 \\
\texttt{sprache} & 10 \\
\texttt{genre} & 10 \\
\texttt{autor} & 12 \\
\texttt{buch} & 12 \\
\texttt{buch\_autor} & 14 \\
\texttt{buch\_genre} & 14 \\
\texttt{buchexemplar} & 15 \\
\texttt{zeitslot} & 15 \\
\texttt{standort} & 10 \\
\texttt{versandoption} & 12 \\
\texttt{benachrichtigung} & 12 \\
\texttt{ausleihe} & 12 \\
\texttt{verfuegbarkeit} & 15 \\
\texttt{bewertung} & 10 \\
}
\newcommand{\mdZeilenTabelleA}{%
\texttt{nutzer} & 12 \\
\texttt{adresse} & 12 \\
\texttt{verlag} & 10 \\
\texttt{sprache} & 10 \\
\texttt{genre} & 10 \\
\texttt{autor} & 12 \\
}
\newcommand{\mdZeilenTabelleB}{%
\texttt{buch} & 12 \\
\texttt{buch\_autor} & 14 \\
\texttt{buch\_genre} & 14 \\
\texttt{buchexemplar} & 15 \\
\texttt{zeitslot} & 15 \\
\texttt{standort} & 10 \\
}
\newcommand{\mdZeilenTabelleC}{%
\texttt{versandoption} & 12 \\
\texttt{benachrichtigung} & 12 \\
\texttt{ausleihe} & 12 \\
\texttt{verfuegbarkeit} & 15 \\
\texttt{bewertung} & 10 \\
}

% automatisch erzeugt von index_benchmark.py -- nicht von Hand bearbeiten
\newcommand{\benchmeta}{736.711 Datensätze, Median über 40 Parameterwerte, SQLite 3.50.4}
\newcommand{\benchsize}{35,4~MB $\rightarrow$ 48,7~MB (+38\,\%)}
\newcommand{\benchinsert}{114~ms $\rightarrow$ 195~ms (+70\,\%)}
\newcommand{\benchrows}{%
Q1 & Mein Regal (Exemplare je Eigentümer:in) & 3,22 & 0,05 & 64,6\,$\times$ \\
Q2 & Meine Ausleihen, neueste zuerst & 8,59 & 0,07 & 118,4\,$\times$ \\
Q3 & Mahnlauf: überfällige Ausleihen (2 \%) & 12,64 & 6,71 & 1,9\,$\times$ \\
Q4 & Gegenprobe: zurückgegebene Ausleihen (90 \%) & 85,81 & 96,55 & 0,9\,$\times$ \\
Q5 & Werk $\rightarrow$ alle Exemplare & 3,11 & 0,06 & 51,8\,$\times$ \\
Q6 & Umkreissuche 15 km (Haversine + View) & 918,04 & 1,63 & 561,7\,$\times$ \\
Q7 & Titelsuche Präfix (LIKE 'Wort\%') & 0,28 & 0,05 & 5,9\,$\times$ \\
Q8 & Badge: ungelesene Benachrichtigungen & 7,91 & 0,03 & 275,5\,$\times$ \\
Q9 & Alle Bücher eines Genres & 0,49 & 0,26 & 1,9\,$\times$ \\
}


\lstdefinestyle{sql}{
  language=SQL, basicstyle=\ttfamily\scriptsize, keywordstyle=\color{blue!60!black}\bfseries,
  commentstyle=\color{gray}, stringstyle=\color{teal!70!black}, columns=fullflexible,
  breaklines=true, frame=single, rulecolor=\color{gray!40}, numbers=left,
  numberstyle=\tiny\color{gray}, upquote=true, morekeywords={PRAGMA,TRIGGER,RAISE,ANALYZE,EXPLAIN,QUERY,PLAN,AUTOINCREMENT,VIEW,WITH,COLLATE,NOCASE,ESCAPE},
  literate={ö}{{\"o}}1 {ä}{{\"a}}1 {ü}{{\"u}}1 {Ö}{{\"O}}1 {Ä}{{\"A}}1 {Ü}{{\"U}}1 {ß}{{\ss}}1
}
\lstdefinestyle{shell}{basicstyle=\ttfamily\small, frame=single, rulecolor=\color{gray!40},
  backgroundcolor=\color{gray!6}, columns=fullflexible, breaklines=true}

% Screenshot aus dem DBMS einbinden; fehlt die Datei, erscheint ein Platzhalter.
% Optionales 1. Argument: maximale Bildhoehe (Standard 7cm; hochkante Bilder brauchen mehr).
\newcommand{\shot}[3][7cm]{%
  \begin{center}
  \IfFileExists{screenshots/#2}{%
    \includegraphics[width=0.92\textwidth,height=#1,keepaspectratio]{screenshots/#2}}{%
    \fbox{\parbox[c][2.2cm][c]{0.88\textwidth}{\centering\small
      Platzhalter: Screenshot \texttt{screenshots/#2} fehlt noch\\(siehe \texttt{screenshots/README.md})}}}\\[2pt]
  {\small\textit{#3}}
  \end{center}}

\title{\vspace{-1.5cm}\Large\textbf{Portfolioteil 3 -- SQL-Dokumentation}\\[2pt]
\large SQL-Dateien, Installationsanleitung und Ergebnisse der Finalisierungsphase\\[2pt]
\normalsize DLBDSPBDM01\_D -- Data-Mart-Erstellung in SQL}
\author{Tizian Senger -- Matrikelnummer: IU14143428\\
\url{https://github.com/TizianSenger/IUDatabaseProject}}
\date{\today}

\begin{document}
\maketitle
\vspace{-0.8cm}
\tableofcontents

% =========================================================
\section{Dateien}
% =========================================================
\begin{center}\small
\begin{tabular}{@{}p{0.47\textwidth}p{0.47\textwidth}@{}}
\toprule
\textbf{Datei (Ordner \texttt{03\_Finalisierungsphase})} & \textbf{Inhalt} \\
\midrule
\texttt{Senger-Tizian\_IU14143428\_DLBDSPBDM01\_D\_\allowbreak Abschlussphase\_SQL.sql} & Finales, eigenständiges SQL-Skript: Schema, Trigger, Dummy-Daten, Testfälle, Indizes, View, Umkreissuche, Metadaten (Anhang~A). \\
\texttt{install\_db.py} & Baut \texttt{buchtausch\_app.db} auf, führt alle Abfragen und die sechs Negativtests aus, schreibt \texttt{metadaten.tex}. \\
\texttt{index\_benchmark.py} & Nutzennachweis der Indizes auf einer separaten Benchmark-Datenbank (Abschnitt~4.3). \\
\texttt{buchtausch\_app.db} & Fertig erzeugte SQLite-Datenbank (\mdKiB{} KiB). \\
\texttt{benchmark\_ergebnisse.md} & Vollständige Benchmark-Ergebnisse inkl. aller Ausführungspläne. \\
\bottomrule
\end{tabular}
\end{center}

% =========================================================
\section{Installationsanleitung}
% =========================================================
\textbf{Voraussetzung:} SQLite ab Version 3.35 (für die Mathe-Funktionen \texttt{sin}, \texttt{cos}, \texttt{asin}, \texttt{radians} der Haversine-Formel). Ein Datenbankserver ist nicht nötig.

\textbf{Variante A -- Python (empfohlen, keine Zusatzinstallation):} Python~3 bringt das Modul \texttt{sqlite3} bereits mit.
\begin{lstlisting}[style=shell]
cd latex/03_Finalisierungsphase
python install_db.py          # DB aufbauen + alle Tests ausgeben
python install_db.py --quiet  # nur Aufbau + Kurzbericht (Negativtests)
python index_benchmark.py     # optional: Index-Benchmark (ca. 1 Minute)
\end{lstlisting}
Das Skript beendet sich mit Exit-Code~0, wenn alle sechs Negativtests wie erwartet scheitern.

\textbf{Variante B -- SQLite-Kommandozeile:}
\begin{lstlisting}[style=shell]
sqlite3 buchtausch_app.db < Senger-Tizian_IU14143428_DLBDSPBDM01_D_Abschlussphase_SQL.sql
\end{lstlisting}

\textbf{Variante C -- grafisch:} Die mitgelieferte \texttt{buchtausch\_app.db} direkt in einem SQLite-Werkzeug öffnen (z.\,B. DB Browser for SQLite oder eine VS-Code-SQLite-Erweiterung) und die Abfragen aus dem Skript einzeln ausführen.

\textbf{Wichtig -- Fremdschlüssel:} SQLite setzt \texttt{FOREIGN KEY}-Constraints nur durch, wenn in \emph{derselben Verbindung} vorher \texttt{PRAGMA foreign\_keys = ON;} ausgeführt wurde. Das Skript und \texttt{install\_db.py} tun das automatisch. In grafischen Werkzeugen, die je Abfrage eine neue Verbindung öffnen, muss das PRAGMA vor dem FK-Negativtest im selben Abfragefenster stehen.

% =========================================================
\section{Aufbau des SQL-Skripts}
% =========================================================
\begin{enumerate}
  \item PRAGMA-Einstellungen
  \item \texttt{CREATE TABLE} für \mdTabellen{} Tabellen mit PK, FK, UNIQUE, CHECK, NOT NULL, DEFAULT
  \item \mdTrigger{} Trigger für tabellenübergreifende Geschäftsregeln
  \item Dummy-Daten (\mdDatensaetze{} Datensätze, mindestens \mdMinZeilen{} je Tabelle)
  \item Testfälle: (a) eine erfolgreiche Abfrage je Entität, (b) sechs Negativtests
  \item \textbf{neu:} \mdIndizes{} Indizes + \texttt{ANALYZE}
  \item \textbf{neu:} View \texttt{v\_verfuegbares\_angebot} und Umkreissuche per Haversine
  \item \textbf{neu:} Testfälle zu Umkreissuche und Index-Nutzung (\texttt{EXPLAIN QUERY PLAN})
  \item \textbf{neu:} Metadaten-Abfragen (Tabellen, Datensätze, Dateigröße)
\end{enumerate}
Abschnitte 1--5 sind inhaltlich unverändert aus Phase~2 übernommen; die Phase-2-Abgabe liegt unverändert in \texttt{02\_Erarbeitungs-Reflexionsphase}.

% =========================================================
\section{Ergebnisse der Finalisierungsphase}
% =========================================================
\subsection{Umkreissuche (Haversine)}
Suchpunkt Düsseldorf Hbf (51,2200\,N / 6,7942\,E), Radius 60\,km, Ergebnis der Kernabfrage aus Abschnitt~7:

\begin{center}\small
\begin{tabular}{@{}lllllr@{}}
\toprule
\textbf{Titel} & \textbf{Zustand} & \textbf{Eigentümer:in} & \textbf{Standort} & \textbf{Versand} & \textbf{Distanz} \\
\midrule
Ein Skandal in Böhmen & gebraucht & Ben Fischer & Nachbarschaftstreff Köln & -- & 33,0\,km \\
Faust & gut & Clara Hoffmann & Nachbarschaftstreff Köln & beides & 33,0\,km \\
1984 & gut & Anna Bauer & Dortmunder Bucheck & post & 57,0\,km \\
\bottomrule
\end{tabular}
\end{center}
Nicht enthalten, obwohl im Radius: Buchtreff Düsseldorf (1,4\,km, Exemplar nicht freigegeben) und Essener Buchbörse (28,3\,km, Exemplar aktuell verliehen). Kontrollabfrage 8a (Distanz aller Standorte): Düsseldorf 1,4 -- Essen 28,3 -- Köln 33,0 -- Dortmund 57,0 -- Frankfurt 181,4 -- Stuttgart 320,9 -- Hamburg 338,3 -- Leipzig 388,4 -- Berlin 476,3 -- München 485,5\,km. Randfall 8b (10\,km um München): 0 Treffer, da beide Exemplare dort nicht freigegeben sind.

\shot{umkreissuche.png}{Screenshot: Ergebnis der Umkreissuche im DBMS}

\subsection{Index-Nutzung (EXPLAIN QUERY PLAN)}
\begin{center}\footnotesize
\begin{tabular}{@{}lll@{}}
\toprule
\textbf{Abfrage} & \textbf{ohne Index} & \textbf{mit Index} \\
\midrule
Q1 Mein Regal & \texttt{SCAN e} & \texttt{SEARCH e USING INDEX idx\_buchexemplar\_eigentuemer} \\
Q2 Meine Ausleihen & \texttt{SCAN} + \texttt{TEMP B-TREE} & \texttt{SEARCH USING INDEX idx\_ausleihe\_nutzer\_datum} \\
Q3 Mahnlauf & \texttt{SCAN ausleihe} & \texttt{SEARCH USING INDEX idx\_ausleihe\_status} \\
Q4 Bounding Box & \texttt{SCAN standort} & \texttt{SEARCH USING COVERING INDEX idx\_standort\_geo} \\
Q5 Titelsuche & \texttt{SCAN buch} & \texttt{SEARCH USING COVERING INDEX idx\_buch\_titel} \\
Q6 Badge & \texttt{SCAN benachrichtigung} & \texttt{SEARCH USING COVERING INDEX idx\_benachrichtigung\_ungelesen} \\
Q7 Genre & \texttt{SCAN bg} & \texttt{SEARCH bg USING INDEX idx\_buch\_genre\_genre} \\
\bottomrule
\end{tabular}
\end{center}
\textit{Hinweis:} Die Nummerierung entspricht Abschnitt~8c des SQL-Skripts; der Benchmark in~4.3 verwendet eine eigene Nummerierung.

\shot{explain.png}{Screenshot: EXPLAIN QUERY PLAN mit Indexnutzung im DBMS}

\subsection{Laufzeit-Benchmark}
Separate Datenbank mit identischem Schema, \benchmeta:
\begin{center}\small
\begin{tabular}{@{}llrrr@{}}
\toprule
\textbf{Nr.} & \textbf{Abfrage} & \textbf{ohne [ms]} & \textbf{mit [ms]} & \textbf{Faktor} \\
\midrule
\benchrows
\bottomrule
\end{tabular}
\end{center}
Kosten: Dateigröße \benchsize; 5.000 Ausleihen + 5.000 Benachrichtigungen einfügen: \benchinsert. Bewertung siehe Abschlussdokument (BE), Abschnitt~2.1.

\subsection{Negativtests}
\begin{center}\small
\begin{tabular}{@{}lll@{}}
\toprule
\textbf{Nr.} & \textbf{Ungültiger Datensatz} & \textbf{Fehlermeldung (erwartet und erhalten)} \\
\midrule
1 & Sternebewertung 7 & \texttt{CHECK constraint failed: sternebewertung BETWEEN 1 AND 5} \\
2 & Versandkosten $-5{,}00$ & \texttt{CHECK constraint failed: kosten >= 0} \\
3 & Zeitslot 15:00--14:00 & \texttt{CHECK constraint failed: startzeit < endzeit} \\
4 & \texttt{nutzer\_id} 9999 & \texttt{FOREIGN KEY constraint failed} \\
5 & Eigenausleihe & \texttt{Nutzer kann eigenes Buchexemplar nicht ausleihen} \\
6 & Bewertung vor Rückgabe & \texttt{Bewertung nur fuer abgeschlossene \dots{} erlaubt} \\
\bottomrule
\end{tabular}
\end{center}

\subsection{Metadaten}
\mdTabellen{} Tabellen, \mdDatensaetze{} Datensätze, \mdIndizes{} eigene und \mdAutoIndizes{} automatische Indizes, \mdTrigger{} Trigger, \mdViews{} View; Dateigröße \mdBytes{} Byte (\mdSeiten{} Seiten à \mdSeitengroesse{} Byte = \mdKiB{} KiB), SQLite \mdSQLite.

\shot[13cm]{metadaten.png}{Screenshot: Metadaten-Abfragen (Abschnitt 9) im DBMS}

% =========================================================
\appendix
\section{Vollständiges SQL-Skript}
% =========================================================
\lstinputlisting[style=sql]{Senger-Tizian_IU14143428_DLBDSPBDM01_D_Abschlussphase_SQL.sql}

\end{document}
